\cdcsectionnonnum{Glossaire et acronymes}

\cdcintertitre{Termes métier}

\begin{cdcdefinitions}
  \item[Qribna] Nom de la plateforme web décrite dans ce document.
  \item[Client] Personne qui recherche une boutique et confie son téléphone à réparer.
  \item[Propriétaire] Personne qui crée et gère une boutique sur la plateforme, ses réparateurs, ses tarifs et ses revenus.
  \item[Réparateur] Personne rattachée à une boutique, qui enregistre les appareils reçus et effectue les réparations.
  \item[Administrateur] Personne chargée de la supervision de la plateforme (utilisateurs, boutiques, avis, signalements).
  \item[Boutique] Point de réparation de téléphones inscrit sur la plateforme, avec une adresse, des horaires et des services.
  \item[État initial] État physique et fonctionnel du téléphone à sa réception, documenté par des photos et validé par le client.
  \item[Numéro de suivi] Identifiant unique attribué à chaque réparation, qui permet au client d’en suivre l’avancement.
  \item[Statut de réparation] Étape courante d’une réparation (reçu, en diagnostic, en réparation, prêt à récupérer, livré).
  \item[Diagnostic préliminaire] Estimation indicative des causes, du prix et du délai, fournie avant l’examen de l’appareil et confirmée ensuite par le réparateur.
  \item[Prestataire de paiement] Organisme tiers qui traite les paiements électroniques pour le compte de la plateforme.
  \item[Commission] Montant calculé selon des règles définies sur les réparations payées.
  \item[Reçu numérique] Justificatif électronique d’un paiement, associé à une réparation.
  \item[Garantie] Conditions et durée pendant lesquelles une réparation est couverte par la boutique.
  \item[Observatoire des prix] Module d’analyse statistique agrégée des prix, des pannes et des délais de réparation au Maroc.
  \item[Darija] Arabe dialectal marocain.
\end{cdcdefinitions}

\cdcintertitre{Acronymes et abréviations}

\begin{cdcdefinitions}[2.2cm]
  \item[CDCF] Cahier des charges fonctionnel.
  \item[CNDP] Commission nationale de contrôle de la protection des données à caractère personnel.
  \item[MAD] Dirham marocain, monnaie des montants affichés.
  \item[MoSCoW] Méthode de priorisation (\prioMust, \prioShould, \prioCould, \prioWont).
  \item[V1, V2] Version~1 (fonctionnalités essentielles) et version~2 (fonctionnalités avancées) de la plateforme.
\end{cdcdefinitions}
